\chapter{Interfaccia host SPI}
\label{ch:spi}

\section{Livello fisico}
L'FPGA è sempre \textbf{slave} SPI. Il protocollo v1 usa SPI \textbf{Mode~0}
(CPOL=0, CPHA=0), MSB-first, single-SPI. Un comando per periodo di CS basso; il byte~0
di ogni transazione è l'opcode. I campi multi-byte sono big-endian.

\begin{fnspec}[Campionamento Mode 0]
\code{mosi} è campionato sul fronte di \textbf{salita} di \code{sclk}; \code{miso} è
pilotato sul fronte di \textbf{discesa} (stabile prima del successivo campionamento del
master). \code{spi\_slave} sincronizza \code{sclk/mosi/cs\_n} con un doppio flip-flop
(CDC a 3 stadi) prima di ogni rilevazione di fronte.
\end{fnspec}

\begin{center}
\begin{tikztimingtable}[timing/dslope=0.1,timing/.style={x=3.4ex,y=2.2ex},
  xscale=1.0,font=\scriptsize]
  \sig{CS\_N}   & H 1L 16L 1H \\
  \sig{SCLK}    & L 1L {2C(2)}8{2C(2)} 6L \\
  \sig{MOSI}    & U 1U 2D{b7} 2D{b6} 2D{b5} 2D{b4} 2D{b3} 2D{b2} 2D{b1} 2D{b0} 2U \\
  \sig{MISO}    & Z 1Z 16D{dato} 1Z \\
\end{tikztimingtable}
\end{center}
\begin{center}\footnotesize\itshape\color{fnGrey}
Framing di un byte: CS scende, 8 colpi di SCLK, MSB per primo; MISO in tri-state fuori
transazione.\end{center}

\begin{fnnote}[Contratto \texttt{tx\_byte\_req}]
\code{tx\_byte\_req} è un \emph{prefetch hint}, non un evento ``byte consumato'': un
consumatore deve avanzare i puntatori (indirizzo RAM, indice byte di risposta) su
\code{rx\_valid}, che pulsa esattamente una volta per byte reale trasferito.
\end{fnnote}

\section{Framing e lunghezza esplicita}
La lunghezza dei trasferimenti RAM è \textbf{esplicita}, non delimitata dal fronte di
CS: \op{WRITE\_RAM}/\op{READ\_RAM} portano un campo lunghezza a 2~byte, così il
controller SPI necessita solo di un contatore di byte. Gli indirizzi di byte sono a
23~bit, trasportati in un campo di 3~byte con il bit più alto riservato a 0.

\section{Tabella degli opcode}
\renewcommand{\arraystretch}{1.16}
\begin{longtable}{C{1.1cm} L{2.4cm} L{3.9cm} L{2.4cm} L{4.0cm}}
\toprule
\rowh \thd{Op} & \thd{Nome} & \thd{Payload (host$\to$FPGA)} & \thd{Risposta} & \thd{Funzione} \\
\midrule
\endfirsthead
\rowh \thd{Op} & \thd{Nome} & \thd{Payload} & \thd{Risposta} & \thd{Funzione} \\ \midrule
\endhead
\bottomrule
\endfoot
\op{0x00} & NOP & --- & --- & Nessuna operazione (idle/dummy clocking). \\
\rowa \op{0x01} & WRITE\_RAM & addr(3B)+len(2B)+dati & --- & Scrive un blocco in PSRAM (X, pesi, bias, parametri). \\
\op{0x02} & READ\_RAM & addr(3B)+len(2B) & \code{len} byte & Rilegge un blocco da PSRAM. \\
\rowa \op{0x0F} & RESET & --- & --- & Reset sincrono del motore e azzeramento del latch STATUS; non cancella la PSRAM. \\
\op{0x10} & SET\_BASE & sel(1B)+addr(3B) & --- & Imposta le basi/registri (vedi §\ref{sec:setbase}). \\
\rowa \op{0x11} & SET\_NET\_TYPE & type(1B) & --- & Tipo di rete: \code{0x01}=dense (\#1), \code{0x02}=graph (\#2). Default dopo RESET=dense. \\
\rowa \op{0x20} & START & --- & --- & Avvia \code{neuron\_memory} (percorso single-layer); ignorato se busy. \\
\op{0x21} & STATUS & --- & 1 byte & bit0=\code{busy} (live), bit1=\code{done} (sticky, clear-on-read), bit2=\code{err} (guard grafo), bit3=\code{flash\_err} (sticky, clear-on-read), bit4=\code{flash\_busy} (live); bit7:5=0. \\
\rowa \op{0x22} & READ\_OUTPUT & --- & \code{N\_NEURONS} byte & \code{y\_bus} neuron-major (byte~0 = neurone~0); solo percorso dense (Tipo \#1). \\
\op{0x23} & RUN\_NETWORK & num\_layers(1B) & --- & Avvia l'esecuzione: dispatch su \code{net\_type} verso \code{layer\_sequencer} (\#1) o \code{graph\_engine} (\#2); ignorato se busy. \\
\rowa \op{0x30} & READ\_CONFIG & --- & 11 byte & Record di configurazione hardware (§\ref{sec:readcfg}). \\
\op{0x40} & FLASH\_READ\_BLOCK & flash\_addr(3B)+psram\_addr(3B)+len(3B) & --- & Lettura raw flash$\to$PSRAM, bypassa il catalogo. \\
\rowa \op{0x41} & FLASH\_WRITE\_BLOCK & psram\_addr(3B)+flash\_addr(3B)+len(3B) & --- & Scrittura raw PSRAM$\to$flash (erase-before-write interno + loop Page Program $\leq$256B + poll WIP, trasparente all'host), bypassa il catalogo. \\
\op{0x42} & FLASH\_ERASE & sector\_addr(3B) & --- & Erase di un settore da 4~KB (deve essere sector-aligned), bypassa il catalogo. \\
\rowa \op{0x43} & CAT\_READ & --- & --- & Ricarica il catalogo a 16 slot (registri on-chip) dal settore riservato in flash. \\
\op{0x44} & CAT\_WRITE\_SLOT & slot\_id(1B)+offset(3B)+len(3B)+tipo(1B) & --- & Registra/aggiorna (offset, lunghezza, tipo) dello slot nel catalogo on-chip e lo persiste in flash; marca lo slot \emph{non valido} finché \op{SAVE\_SLOT} non lo conferma. \\
\rowa \op{0x45} & LOAD\_SLOT & slot\_id(1B)+psram\_addr(3B) & --- & Flash$\to$PSRAM per lo slot (offset/lunghezza dal catalogo), verifica CRC32 live; \code{STATUS.flash\_err} se lo slot non è valido o il CRC non torna. \\
\op{0x46} & SAVE\_SLOT & slot\_id(1B)+psram\_addr(3B)+len(3B) & --- & PSRAM$\to$flash all'offset già registrato dello slot, calcola il CRC32 live; a esito positivo aggiorna e persiste la entry di catalogo (lunghezza, CRC, valid=1). \\
\rowa \op{0x47} & CAT\_INSPECT & slot\_id(1B) & 16 byte & Lettura sincrona di una entry di catalogo già caricata: offset[3]+len[3]+tipo[1]+valid[1]+CRC32[4]+riservato[4], MSB-first. \\
\end{longtable}
Tutti gli opcode flash sono \emph{fire-and-forget}: l'host fa polling su \op{STATUS}
(bit4=\code{flash\_busy}, bit3=\code{flash\_err}) o sui pin \code{irq\_n}/\code{data\_ready\_n}
per l'esito, eccetto \op{CAT\_INSPECT} che risponde in modo sincrono.

Gli 8 opcode flash (\op{0x40}--\op{0x47}) sono descritti in dettaglio, con
razionale di progetto e latenze reali misurate, in §\ref{sec:flashspi} sotto.

\section{Selettori \texttt{SET\_BASE}}
\label{sec:setbase}
\begin{tabularx}{\textwidth}{C{1.2cm} L{3.2cm} Y}
\toprule
\rowh \thd{sel} & \thd{Registro} & \thd{Uso} \\
\midrule
0 & \code{x\_base} & Base ingresso $X$. \\
\rowa 1 & \code{w\_base} & Base pesi. \\
2 & \code{bias\_addr} & Base bias. \\
\rowa 3 & \code{table\_base} & Base tabella descrittori (multi-layer). \\
4 & \code{buf\_a\_base} & Buffer ping-pong A. \\
\rowa 5 & \code{buf\_b\_base} & Buffer ping-pong B. \\
6 & \code{activation} & Attivazione (2 bit bassi) --- solo percorso single-layer. \\
\rowa 7 & \code{n\_inputs\_real} & Larghezza ingressi runtime (16-bit BE) --- single-layer. \\
8 & \code{n\_neurons\_real} & Larghezza neuroni runtime (16-bit BE) --- single-layer. \\
\rowa 9 & \code{num\_neurons\_graph} & Numero neuroni del grafo (16-bit BE) --- Tipo \#2. \\
10 & \code{n\_out} & Numero id di uscita (16-bit BE) --- Tipo \#2. \\
\bottomrule
\end{tabularx}
I selettori 6--8 riguardano solo il percorso single-layer/manuale; con \op{RUN\_NETWORK}
i valori equivalenti sono letti per-layer dalla tabella descrittori.

\begin{fnwarn}[Casi limite ``reale=0'' corretti (2026-09-04)]
La campagna di ri-certificazione (\code{docs/validation/bugs.md}) ha trovato che diversi
valori runtime pari a zero non erano protetti da alcun guard, con esiti che andavano da un
risultato silenziosamente ignorato fino a hang o scritture PSRAM a indirizzi arbitrari.
Tutti e cinque i casi seguenti sono ora no-op sicuri, verificati indipendentemente:
\begin{itemize}
\item \code{n\_inputs\_real=0} (selettore 7): completa in 1 ciclo con
$y=\text{activation}(\text{bias})$ (BUG-003).
\item \code{n\_neurons\_real=0} (selettore 8): completa senza eseguire alcun calcolo
per-neurone, molto più rapido di un run a piena larghezza (BUG-004).
\item \code{num\_neurons\_graph=0} (selettore 9): completa immediatamente dopo la copia
degli ingressi, senza mai entrare nel loop dei descrittori (BUG-006).
\item \op{RUN\_NETWORK} con \code{num\_layers=0} (percorso dense): no-op immediato ---
\textbf{prima del fix eseguiva 256 layer fasulli leggendo dati PSRAM arbitrari come
descrittori} (BUG-005, CRITICO, vedi \S\ref{sec:run-network} sotto).
\item \op{SET\_NET\_TYPE} ricevuto mentre un run è in corso: ora rifiutato silenziosamente
(nessun effetto, nessun errore SPI) invece di rimappare il multiplexer dell'arbitro a metà
esecuzione --- \textbf{prima del fix causava un hang permanente del motore in corso}
(BUG-007, CRITICO).
\end{itemize}
Dettagli, evidenza e verifica di ciascun fix in \code{docs/validation/bugs.md}.
\end{fnwarn}

\section{\texttt{STATUS.done} sticky / clear-on-read}
In \code{neuron\_memory} il segnale \code{done} è un impulso di un solo ciclo. Un host
che effettua polling via SPI (molto più lento del clock FPGA) mancherebbe quasi
certamente un impulso grezzo di un ciclo. Il banco registri SPI latcha quindi
\code{done} in un bit sticky sull'impulso e lo azzera quando l'host legge \op{STATUS}
(o \op{RESET}). Il bit \code{busy} è invece mantenuto a livello per tutta la
computazione e si legge live.

\begin{fnwarn}[Race corretto (2026-09-02)]
Una race reale nel meccanismo sticky (presente dalla Fase~4) è stata corretta latchando
uno \code{status\_snapshot} all'accettazione dell'opcode \op{STATUS} e condizionando la
pulizia del bit sticky a \code{status\_snapshot[1]} (si azzera solo se il byte
effettivamente trasmesso mostrava \code{done=1}). Un \code{done} che arriva troppo tardi
per uno snapshot viene riportato al polling successivo invece di essere perso.
\end{fnwarn}

\section{Pin di attenzione host (\texttt{data\_ready\_n}, \texttt{irq\_n})}
Oltre al polling di \op{STATUS}, il top-level espone due pin fisici attivi bassi (banco 7,
cap.~\ref{ch:hw}) che rispecchiano i bit sticky senza richiedere una transazione SPI,
utili per pilotare un GPIO/IRQ dell'host:
\begin{itemize}
\item \code{data\_ready\_n} = $\sim$\code{STATUS.done} (sticky): basso quando un risultato è
pronto da leggere, torna alto alla lettura di \op{STATUS} (clear-on-read).
\item \code{irq\_n} = $\sim$\code{STATUS.err} (guard grafo): basso quando il guard load-time
di \code{graph\_engine} è scattato. \textbf{Non} è clear-on-read: si azzera solo con
\op{RESET} o un nuovo avvio di grafo, così un errore non passa inosservato tra un polling e
l'altro.
\end{itemize}
Sono porte aggiuntive: non toccano gli opcode né i registri esistenti.

\begin{fnwarn}[\code{flash\_err} non ha un pin dedicato]
\code{STATUS.flash\_err} (bit3) è riportato \textbf{solo} nel byte \op{STATUS}, per scelta
di progetto: riusare \code{irq\_n} lo avrebbe confuso con gli errori del guard grafo (due
domini di errore indipendenti sullo stesso pin), mentre un'operazione flash è sempre
avviata dall'host con un opcode appena emesso, quindi il polling di \op{STATUS} subito dopo
--- già implicito nella convenzione ``fire-and-forget, poi polling \op{STATUS}/
\code{data\_ready\_n}'' --- è già naturale, senza bisogno di un pin asincrono in più.
\code{data\_ready\_n} invece \emph{si azzera anche al termine di un'operazione flash}: lo
specchia \code{STATUS.done} (bit1), che ora latcha anche sul completamento di un op flash,
non solo su \op{RUN\_NETWORK}/\op{START}.
\end{fnwarn}

\section{\texttt{READ\_CONFIG}}
\label{sec:readcfg}
Payload fisso di \textbf{11 byte}: permette a un unico firmware host di funzionare con
bitstream diversi senza ricompilare. I valori \code{N\_INPUTS}/\code{N\_NEURONS} riportano
il \emph{massimo} di build (il soffitto), non necessariamente la rete correntemente
caricata.

\begin{tabularx}{\textwidth}{C{1.6cm} L{3.6cm} Y}
\toprule
\rowh \thd{Byte} & \thd{Campo} & \thd{Sorgente} \\
\midrule
0 & \code{ADDR\_WIDTH} (bit) & \code{neuron\_memory.ADDR\_WIDTH} \\
\rowa 1--2 & \code{N\_INPUTS} (16-bit BE) & massimo di build \\
3 & \code{N\_NEURONS} & massimo di build \\
\rowa 4 & \code{PARALLEL} & parametro di build \\
5 & \code{DATA\_WIDTH} (bit) & parametro di build \\
\rowa 6--7 & versione protocollo (BE) & \code{0x0001} \\
8--9 & \code{N\_TOTAL} (16-bit BE) & massimo segnali grafo (Tipo \#2) \\
\rowa 10 & flag di capacità & bit0=\code{GRAPH\_SUPPORTED}=1 \\
\bottomrule
\end{tabularx}

\section{Sottosistema flash (opcode 0x40--0x47, completato 2026-09-04)}
\label{sec:flashspi}
La FPGA ha accesso \textbf{esclusivo} alla flash di boot/persistenza onboard (Winbond
\code{W25Q128JV}, 16~MB SPI NOR, cap.~\ref{ch:hw} §6/§7) tramite un SPI master dedicato e
fisicamente separato (\code{rtl/spi\_flash\_master.v}), mai per accesso diretto dell'host ai
pin della flash. \textbf{Non} è un filesystem: un catalogo a dimensione fissa (16 slot,
\code{rtl/flash\_slot\_manager.v}) mappa \code{slot\_id}~$\to$~(offset, lunghezza, tipo,
valid, CRC32) in un settore riservato della flash (settore 0) --- nessuna allocazione
dinamica, nessun garbage collection.

\begin{fnnote}[Stratificazione (ogni livello testabile a sé)]
\begin{itemize}
\item \code{rtl/spi\_flash\_master.v} --- SPI master grezzo verso il chip flash
      (RDID/READ/WREN/PP/SE/RDSR-1). Bus a 4 fili completamente indipendente
      (\code{sclk}/\code{mosi}/\code{miso}/\code{cs\_n}, tutti GPIO ordinario ---
      Fase F7, 2026-09-04): una versione precedente riusava il pad \code{CCLK} di boot
      via la primitiva ECP5 \code{USRMCLK} per risparmiare un pin, abbandonato perché
      rendeva fuorviante l'affermazione di ``bus esclusivo'' (elettricamente dipendeva
      comunque dal motore di configurazione) e comportava un gap di verifica mai chiuso
      (timing di \code{USRMCLKTS} mai verificato contro la guida Lattice primaria).
\item \code{rtl/flash\_copy\_engine.v} --- motore di streaming a blocchi: flash$\to$PSRAM
      (\code{DIR\_LOAD}), PSRAM$\to$flash con erase-before-write interno + loop Page
      Program $\leq$256B + poll WIP (\code{DIR\_SAVE}), erase di settore standalone
      (\code{DIR\_ERASE}). Master a bassa priorità (Porta D) su \code{rtl/mem\_arbiter.v}:
      le operazioni flash sono su scala dei ms e non bloccano mai l'inferenza.
\item \code{rtl/flash\_slot\_manager.v} --- il catalogo a slot sopra, più un CRC32
      (\code{rtl/crc32.v}, IEEE~802.3/zlib) calcolato live sul flusso di byte reale durante
      \op{LOAD\_SLOT}/\op{SAVE\_SLOT}, così uno slot corrotto o scritto a metà (es.
      alimentazione persa durante l'erase) è rilevato anche quando l'operazione flash
      sottostante ha riportato successo.
\end{itemize}
\end{fnnote}

\begin{fnwarn}[Allineamento a settore obbligatorio]
\op{SAVE\_SLOT} (e i raw \op{FLASH\_WRITE\_BLOCK}/\op{FLASH\_ERASE}) richiedono che
l'indirizzo flash target sia allineato a settore da 4~KB --- rifiutato come errore
altrimenti, invece di un silenzioso read-modify-erase-write parziale del settore (non
esiste un buffer di scratch abbastanza grande per farlo, e ogni \op{SAVE\_SLOT} reale scrive
già uno slot intero e allineato per costruzione).
\end{fnwarn}

Razionale completo, ogni citazione da datasheet, ogni test avversariale (CRC non
corrispondente, slot mai salvato, attraversamento di confine pagina, simulazione di perdita
di alimentazione, contesa sull'arbitro) e i due bug reali trovati e corretti durante il
bring-up (uno pre-esistente in \code{psram\_controller.v}, uno nel nuovo handshake di
richiesta dell'arbitro) sono in \code{WORKLOG.md} (voci Fasi F1-F6) e
\code{docs/FPGA-Neural-Flash-Subsystem-Verification.md} (sunto di copertura per modulo, non
ripetuto qui).

\begin{tabularx}{\textwidth}{L{3.4cm}Y}
\toprule
\rowh \thd{Operazione} & \thd{Latenza reale misurata} \\
\midrule
ERASE (settore 4~KB) & $\approx$400~ms (dominata dal tSE interno del chip flash, indipendente dal clock host) \\
\rowa SAVE (pagina 256~B, incl. erase interno) & $\approx$403~ms (idem, tSE+tPP) \\
LOAD (4096~B) & 1.74~ms (2.35~MB/s) @80~MHz; 8.71~ms (0.47~MB/s) @16~MHz (solo SPI-clock-bound) \\
\bottomrule
\end{tabularx}
Metodologia di misura completa in \code{docs/FPGA-Neural-Flash-Subsystem-Verification.md}.

\section{Sequenze di sessione}
\subsection{Percorso single-layer}
\begin{lstlisting}[language=,caption={Sessione single-layer},basicstyle=\ttfamily\scriptsize]
RESET               -> 0x0F
READ_CONFIG         -> 0x30      (l'host apprende N_INPUTS/N_NEURONS/...)
WRITE_RAM (pesi)    -> 0x01 ...
WRITE_RAM (bias)    -> 0x01 ...
SET_BASE (X/W/BIAS) -> 0x10 x3
WRITE_RAM (input X) -> 0x01 ...
START               -> 0x20
poll STATUS         -> 0x21      (finche' done=1; si azzera a questa lettura)
READ_OUTPUT         -> 0x22
\end{lstlisting}

\subsection{Percorso multi-layer (RUN\_NETWORK)}
\label{sec:run-network}
\begin{lstlisting}[language=,caption={Sessione multi-layer},basicstyle=\ttfamily\scriptsize]
WRITE_RAM (tabella descrittori)          -> 0x01 ...
WRITE_RAM (pesi/bias per layer, X layer0)-> 0x01 ...
SET_BASE (X/TABLE/BUF_A/BUF_B)           -> 0x10 x4
RUN_NETWORK(num_layers)                  -> 0x23 <num_layers>
poll STATUS                              -> 0x21   (finche' done=1)
READ_OUTPUT                              -> 0x22   (y_bus del layer finale)
\end{lstlisting}

\begin{fnnote}[Fuori ambito per v1]
Dual~SPI e CRC/checksum sui trasferimenti host (SPI assunto affidabile su traccia di
scheda --- da non confondere con il CRC32 del catalogo flash, §\ref{sec:flashspi}, che
protegge un dominio diverso: la persistenza flash$\leftrightarrow$PSRAM, non il link SPI
host).
\end{fnnote}

\begin{fnwarn}[\op{WRITE\_RAM}/\op{READ\_RAM} senza backpressure verso l'host --- rischio reale, non teorico]
Ogni byte ricevuto/prodotto deve essere completamente processato da \code{spi\_engine}
prima che arrivi il successivo confine di byte scandito da SCLK --- ragionevole per il
bulk-loading iniziale di pesi/ingressi, non un percorso real-time. Il rischio concreto: se
un host emette \op{WRITE\_RAM}/\op{READ\_RAM} prima che la sequenza di power-up di
\code{psram\_controller.v} sia completata ($\sim$150~\textmu s dopo il reset,
\code{STATE\_INIT}+\code{STATE\_CR\_INIT}), \code{spi\_engine} si blocca in attesa che il
primo accesso PSRAM completi, mentre l'host --- non rallentato da alcun handshake ---
continua a scandire byte. I byte ricevuti durante quello stallo vengono \textbf{scartati
silenziosamente}, senza errore e senza hang: solo dati sbagliati in PSRAM. Trovato durante
il lavoro sul sottosistema flash (\code{WORKLOG.md}, Fase~F5) con una riproduzione minimale
solo-\op{WRITE\_RAM}, senza alcun opcode flash coinvolto: è un rischio generale per
qualunque host, non specifico agli opcode flash. \textbf{Mitigazione attuale: l'host deve
attendere il power-up della PSRAM (o assicurarsi che la FPGA sia fuori reset da
$>$150~\textmu s) prima del suo primo \op{WRITE\_RAM}/\op{READ\_RAM}.} Non risolto a livello
di protocollo (richiederebbe una vera backpressure, una modifica più ampia) ---
dichiarato qui come rischio aperto, non aggirato silenziosamente.
\end{fnwarn}
